\chapter{Rainforest Generators and Terra Preta Rain}
\label{ch:19}

The second terrestrial program concerns the support of biological production on land, and it is stated in the author's own vocabulary as a pair of systems that are designed to operate together. A rainforest generator, in this vocabulary, is a support architecture: a tensile trellis or hanging cable grid carried on smart anchors, with light wells, misters, nutrient fog, and microclimate control, whose function is to enlarge the physical state space available to rainforest growth. It is not a machine that produces a forest from nothing, and the distinction is central to its assessment. A forest is an ecological community whose composition and persistence depend on seed sources, soil organisms, pollinators, and disturbance regimes that no structure supplies, and the generator is intended to supply only what the structure can supply, namely vertical surface on which epiphytes, lianas, and climbers can establish, a controlled supply of moisture and nutrients to the canopy, and a modified light and temperature regime beneath it. Terra preta rain is the paired system for the supply of water and nutrients. Rainwater is collected in elevated basins, including the basins of calderas, filtered through gravel, biochar, and sand and through biochar-lined berms, enriched with nutrients, minerals, microbes, and suspended carbon, and redistributed under gravity over the constructed canopy as a circulating medium and not a static amendment. The name recalls the Amazonian dark earths, and it signals a hypothesis that the properties of those soils can be reproduced in a flowing form. The agronomic evidence on charcoal-amended soils is reviewed in the literature on biochar~\cite{lehmann2006}, and the extension of those results to a flowing, atmospheric form is the book's hypothesis and is not established by that literature.

The precedent deserves a measured account. The dark earths of the Amazon basin, known by the Portuguese name terra preta de \'indio~\cite{glaser2001}, are patches of soil, often of a few hectares and occasionally larger, that contain high concentrations of charcoal, ceramic fragments, organic matter, and nutrients, and that remain markedly more fertile than the surrounding weathered soils centuries after their formation. They are generally attributed to the deliberate and incidental activities of pre-Columbian populations, who deposited the residues of cooking, waste, and burning, and the persistence of their carbon, which has been attributed to the stability of charred material and to the community of microorganisms associated with it, has made them a model in the literature on soil carbon and on sustainable agriculture in the humid tropics. Several features of the precedent must be noted. The dark earths accumulated over generations and not in a season, which suggests that the process at issue is slow. They are a soil amendment, a body of material in place, and the flowing form proposed here is a different thing whose behavior has to be demonstrated and not inferred. And the mechanisms by which the earths sustain fertility are not fully settled, so that an engineering system modelled on them inherits the uncertainty. The book therefore treats the extension from static soil to circulating medium as a hypothesis for which the program must supply its own evidence.

Three criteria govern the assessment, and they follow from the principle stated in the author's earlier work that the quantity of biomass is not a measure of the health of the biosphere. The first is habitat complexity: the structure should support a diversity of niches and the organisms that occupy them, not merely a large mass of foliage, and a trellis that carried a monoculture of fast-growing climbers would increase biomass while reducing the quality of the habitat. The second is nutrient continuity: the system should maintain a supply of nutrients matched to the demand of the community and should not produce pulses that favor a few opportunistic species or that leave the system through drainage. The third is ecological autonomy: the community that develops should become able to sustain itself without the support of the structure, so that the structure is scaffolding in the literal sense, needed during establishment and removable thereafter. The third criterion has a consequence for the evaluation of persistence which the book adopts as a general principle for interventions of this kind. The risk of withdrawing the support is not a fixed quantity. It is proportional to the fraction of the supported biomass that has not yet become autonomous, and that fraction should be tracked continuously and not asserted once. A system that reports its fraction of autonomous biomass at stated intervals, and that publishes the figures, supplies the evidence on which a decision to continue or to terminate can rationally rest.

The physics of the trellis is that of a hanging cable under load, and a calculation shows that the structural demand is moderate for spans of the order of a hundred meters. A cable grid with parallel cables five meters apart, carrying a distributed load of twenty kilograms per square meter of epiphytes, lianas, water, and accumulated debris, loads each cable with about $1$ kN per meter. Over a span of one hundred meters with a sag of ten meters the horizontal tension is $123$ kN and the tension at the anchors is $132$ kN, which requires a steel cross-section of about $2.6\ \mathrm{cm^2}$ at a working stress of $500$ MPa, a wire rope of roughly eighteen millimeters in diameter. These magnitudes are well within the practice of cable-supported structures. The assumed load is a parameter of the design and not an established figure, since the mass that a constructed canopy would eventually carry depends on what grows on it; mature tropical forest holds on the order of tens of kilograms of above-ground biomass per square meter, of which only a fraction would be borne by a trellis, and the allowances for wind, for the weight of saturated vegetation, and for the dynamic loading of storms would raise the design load considerably. The structural problem is therefore not the difficult one. The difficult problems are ecological and institutional, and the structural calculation shows only that the scale of the structure does not by itself preclude the concept.

The nutrient dynamics of the circulating system have a property that deserves emphasis, because it distinguishes a design that conserves nutrients from one that leaks them. Water delivered to the canopy gives up a fraction of its nutrient content to uptake on each pass and loses another fraction to drainage and leaching, and the remainder returns to the basin. At steady state the fraction of the makeup input that reaches biomass is the ratio of the uptake fraction to the sum of the uptake and loss fractions, and this ratio, which is independent of the rate of circulation, is the efficiency of the loop. If uptake per pass is five per cent and losses are one per cent, the efficiency is $83$ per cent and a nutrient atom makes about seventeen passes before it is lost; if losses rise to five per cent the efficiency falls to $50$ per cent and the number of passes to ten. The result implies that the engineering effort should be directed at reducing loss, by the design of the biochar-lined berms and the capture of drainage, more than at increasing the circulation, and it implies that the monitoring of losses at the boundary of the system is the essential measurement. A system that cannot show that it retains its nutrients is exporting them to the watershed, and the effect on downstream waters could be harmful.

The remaining questions are those of siting and consent. Rainforests grow where rain falls, and the regions suited to the concept are largely in the tropics of Latin America, Africa, and Southeast Asia, much of which is the territory of indigenous and local communities with customary claims, often unrecognized in law. A program that proposed to construct large structures on such land without the free, prior, and informed consent of the communities would reproduce the pattern of extraction against which the book defines itself, and the placement of the structures would be a question of equity and not merely of ecology. The communities hold knowledge of the forests, including in some regions of the dark earths themselves, that no external designer possesses, and the Ostrom principles of Chapter 18 apply here as they did to water: monitoring by the users, graduated sanctions, and recognition of their rights to organize. It is consistent with the program that a community should decline, and a program that could proceed only if no community declined would not be a program of direction.

The assessment under the standard of Chapter 16 places the trellis in the first readiness class in its structural aspects and the system as a whole in the second or third. The mechanism is explicit. The energy balance is simple, since the circulation is gravity-driven and the principal energy cost is the lift required to return water to the basin where rain does not. The material limits are known for the cables and anchors and unknown for the biochar filtration at the scale and flow rates proposed. The comparison with alternatives is required and unflattering to the concept in the near term: the protection of existing forest, the restoration of degraded land by assisted natural regeneration, and the support of indigenous land management all have a better-documented record and lower cost per hectare. The generator and the rain are candidates for supplementing these in places where they are not sufficient, and for demonstrating, at pilot scale and under independent monitoring, that the hypotheses on which they rest are true.

\section*{Formalization}

Four results are given, covering resource limitation, the structure, nutrient retention, and graded persistence.

Let a community grow at a rate determined by resources $R_1,\dots,R_m$ (water, light, nitrogen, phosphorus, and so on) each with requirement $K_i>0$, and let $a_i=R_i/K_i$ be the availability ratio. The \emph{minimum law} states that the relative growth rate is $g=\min\{1,a_1,\dots,a_m\}$.

\begin{proposition}
Under the minimum law, $\partial g/\partial R_i=0$ whenever $a_i>\min_{j}a_j$, so an increase in a non-limiting resource leaves growth unchanged. The marginal value of a resource is positive only for the limiting resources, and if $S$ is the set of indices attaining the minimum with $\min a_j<1$, raising all of $S$ together by a common factor $\lambda$ raises growth to $\min\{1,\lambda\min_ja_j,\min_{i\notin S}a_i\}$.
\end{proposition}

\begin{proof}
The function $\min$ is nondecreasing in each argument and, at a point where the minimum is attained only by index $j$, independent of every other argument locally. The factor formula follows by substituting $\lambda a_i$ for $a_i$, $i\in S$.
\end{proof}

For availability ratios $(N,P,\text{water},\text{light})=(1.4,0.6,1.1,2.0)$ growth is $0.6$ and is limited by phosphorus; doubling nitrogen changes nothing; raising phosphorus to $0.9$ gives growth $0.9$; raising it again to $1.1$ gives growth $1.0$. The implication for design is that the nutrient supply must be matched to the limiting element, which will change as each is relieved, and that a program supplying a standard nutrient mix is wasteful in all but one element.

Second, for a cable of horizontal span $L$ carrying a uniform load $w$ per unit horizontal length with sag $d$, the parabolic approximation gives horizontal tension $H=wL^2/(8d)$, tension at the supports $T_{\max}=\sqrt{H^2+(wL/2)^2}$, and cable length $\approx L\bigl[1+\tfrac83(d/L)^2\bigr]$. The required cross-section for working stress $\sigma$ is $A=T_{\max}/\sigma$. With $w=981$ N m$^{-1}$, $L=100$ m, $d=10$ m: $H=122.6$ kN, $T_{\max}=132.1$ kN, and at $\sigma=500$ MPa, $A=2.64\ \mathrm{cm^2}$ and diameter $18.3$ mm for a solid section of that area. Since $H\propto L^2/d$, halving the sag doubles the tension, which is why the sag is a primary design choice.

Third, consider a circulating loop with basin content $M$ of a nutrient, circulation flow $Q$ of water of concentration $c=M/V$, uptake fraction $f$, loss fraction $\ell$ per pass, and external makeup $m$ per unit time. Then $\dot M=m-(f+\ell)Qc$.

\begin{proposition}
The steady state is $M^\ast=mV/((f+\ell)Q)$ and is globally asymptotically stable, with time constant $\tau=V/((f+\ell)Q)$. The steady uptake is $fm/(f+\ell)$, so the retention efficiency is $\epsilon=f/(f+\ell)$, independent of $Q$ and $V$, and the mean number of passes before loss is $1/(f+\ell)$.
\end{proposition}

\begin{proof}
The equation is linear: $\dot M=m-\frac{(f+\ell)Q}{V}M$, whose solution relaxes exponentially to $M^\ast$ with the stated time constant. Uptake is $fQc^\ast=fm/(f+\ell)$. The number of passes before a unit of nutrient is removed is geometric with parameter $f+\ell$.
\end{proof}

With $f=0.05$ and $\ell=0.01$ one has $\epsilon=0.833$ and $16.7$ passes, and with $\ell=0.05$ one has $\epsilon=0.5$ and $10$ passes. The energy to return water over a head $h$ is $\rho gh$ per unit volume, which is $0.82$ kWh m$^{-3}$ for $h=300$ m, an amount that is avoided entirely where rain refills the basin.

Fourth, let $B$ be the supported biomass and $A(t)\in[0,1]$ the autonomous fraction. If support were withdrawn at time $t$ the expected loss is $L(t)=(1-A(t))B$, the \emph{termination exposure}. If autonomy grows logistically, $A(t)=[1+e^{-k(t-t_0)}]^{-1}$, then the exposure falls below a tolerance $\varepsilon B$ at $t_\varepsilon=t_0+k^{-1}\ln[(1-\varepsilon)/\varepsilon]$. For $k=0.3\ \mathrm{yr^{-1}}$, $t_0=10$ yr and $\varepsilon=0.05$ this gives $t_\varepsilon\approx19.8$ yr. The functional form is a modelling choice, and the point of the statement is that the date at which the structure can be removed is an output of the monitoring, not an input to the design.
